\section{HexStrike Internals and AI Integration}

This chapter focuses on HexStrike: its value, internal architecture, and operational details. HexStrike is the orchestration and automation engine that powers tool selection, execution, and intelligent decision making for the project. The chapter highlights how HexStrike complements the CyberScope pipeline, the internals of agent interaction, orchestration, data processing and the AI integration used for risk assessment.

\section{Why HexStrike matters (Value Proposition)}

HexStrike provides a centralized MCP-compatible orchestration layer that coordinates multiple security tools, optimizes parameters using AI, and aggregates results into actionable intelligence. Key benefits are summarized in Table~\ref{tab:feature-benefits} and visualized in Figure~\ref{fig:feature-box}.

\begin{table}[ht]
\centering
\caption{Feature / Benefit summary}
\label{tab:feature-benefits}
\begin{tabularx}{\textwidth}{|l|X|c|}
\hline
Feature & Problem solved & Impact \\
\hline
Automated orchestration & Eliminates manual coordination across tools and brittle scripts & Repeatable runs \\
\hline
AI-assisted selection & Reduces noisy results and manual parameter tuning & Higher precision \\
\hline
Worker scheduling & Mitigates resource contention across workers & Improved throughput \\
\hline
Caching & Avoids re-running identical scans & Lower cost \\
\hline
\end{tabularx}
\end{table}

\begin{figure}[ht]
\centering
\fbox{\parbox[c][4cm][c]{0.85\textwidth}{\centering System component diagram (placeholder) -- replace with real graphic in \texttt{uploads/diagrams/}}}
\caption{High-level component view (placeholder).}
\label{fig:feature-box}
\end{figure}

% Diagrams removed from this chapter; generated visuals will be inserted during final assembly.

\section{Architecture \& Components}

This section lists the main HexStrike components and responsibilities. See Table~\ref{tab:components} for a compact component summary and Figure~\ref{fig:components} for the system diagram.

\begin{table}[ht]
\centering
\caption{Component summary}
\label{tab:components}
\begin{tabularx}{\textwidth}{|l|X|c|}
\hline
Component & Responsibility & Interfaces / Outputs \\
\hline
MCP Server & Task distribution, agent registry & HTTP / gRPC \\
\hline
AI Decision Engine & Tool selection and parameter optimization & Decision API \\
\hline
Worker Manager & Spawn and schedule workers; health monitoring & Worker queue, health API \\
\hline
Cache & Result deduplication and TTL policies & Key-value store \\
\hline
Database & Persist scan results and queries & SQL / REST \\
\hline
UI Connector & Frontend API surface & REST / WebSocket \\
\hline
\end{tabularx}
\end{table}

\begin{figure}[ht]
\centering
\fbox{\parbox[c][4cm][c]{0.85\textwidth}{\centering Component diagram (placeholder) -- replace with exported diagram}}
\caption{HexStrike component diagram (placeholder).}
\label{fig:components}
\end{figure}

\section{Agent Interaction \& MCP Protocol}

Agents connect to HexStrike using the MCP protocol. They register capabilities, receive tasks, and stream results back. A swimlane-style sequence (Agent / MCP / Worker / Tool) captures registration, task assignment, execution and streaming of results; include an illustrated diagram only if needed.

% Sequence diagram removed: kept concise textual description above.

\section{Tool Orchestration \& Execution Flow}

Tool orchestration covers discovery, selection by the decision engine, execution scheduling, retries and fallback strategies. The text summarizes decision points and parallel execution branches; include a flowchart only when a visual is required for reviewers.

% Orchestration flowchart removed to reduce figures; see Table~\ref{tab:components} and the textual description.

\section{Data Pipeline: Normalization, Matching \& Enrichment}

Raw scanner outputs are transformed into structured \texttt{ScanResult} objects, matched against a vulnerability catalog, and enriched with CVE and CVSS data. Table~\ref{tab:schema} gives a short schema fragment and Figure~\ref{fig:dataflow} shows the data flow.

\begin{table}[ht]
\centering
\caption{Key fields in \texttt{ScanResult} (example)}
\label{tab:schema}
\begin{tabularx}{\textwidth}{|l|l|X|}
\hline
Field & Type & Description \\
\hline
id & string & Unique identifier for the finding \\
\hline
source & string & Tool that produced the finding \\
\hline
evidence & object & Raw evidence and context (headers, request/response) \\
\hline
severity & string & Normalized severity (low/medium/high) \\
\hline
references & array & CVE ids and external links \\
\hline
\end{tabularx}
\end{table}

\begin{figure}[ht]
\centering
\fbox{\parbox[c][3.5cm][c]{0.85\textwidth}{\centering Data pipeline / ER diagram (placeholder)}}
\caption{Data pipeline and entity relationships (placeholder).}
\label{fig:dataflow}
\end{figure}

\section{Process Management, Caching \& Error Recovery}

HexStrike employs process-level state management for workers and a caching layer to avoid re-running expensive scans. Error recovery strategies include exponential backoff, automatic retries, and fallback scanning behaviors. Table~\ref{tab:retry} summarizes retry/backoff policies; include a worker state diagram only if visual clarification is needed.

\begin{table}[ht]
\centering
\caption{Retry/backoff policy (example)}
\label{tab:retry}
\begin{tabular}{|l|l|l|}
\hline
Condition & Action & Notes \\
\hline
Transient network error & Retry up to 3 times & Exponential backoff: 2s, 4s, 8s \\
\hline
Tool timeout & Retry once with increased timeout & Log and escalate after failure \\
\hline
Permanent error & Mark failed & Store diagnostics for analyst \\
\hline
\end{tabular}
\end{table}

% Worker state machine figure removed; see Table~\ref{tab:retry} for policy summary.

\section{Conclusion}

HexStrike acts as the project's orchestration backbone. The architecture and internal workflows presented in this chapter demonstrate how the system automates complex multi-tool assessments and provides higher-level intelligence to reduce analyst workload.